\documentclass[10pt,a4paper]{amsart}
\usepackage[margin=19mm]{geometry}
\usepackage{amsmath,amssymb,mathtools}
\usepackage[scr=rsfs,cal=euler]{mathalfa}
\usepackage{xfrac}
\usepackage[unicode,hidelinks,bookmarks=false]{hyperref}
\newcommand{\ie}{i.e.\ }
\newcommand{\cf}{cf.\ }
\newcommand{\resp}{respectively\ }
\newcommand{\Subsection}{\subsection}
\providecommand{\nobreakdash}{\nobreak\hbox{-}}
\setlength{\emergencystretch}{2em}
\multlinegap0pt
\usepackage{dsfont}
\usepackage{amssymb}
\usepackage{mathtools}
\usepackage{tikz}
\usepackage{bm}
\newtheorem{lemma}{Lemma}
\newcommand{\GL}{\mathrm{GL}}
\newcommand{\mirp}[1]{K_p^{(4)}(#1)}
\title{Eisenstein cohomology and congruences for the ratios of Rankin--Selberg $L$-functions}
\author{P. Narayanan, A. Raghuram}
\date{Source: arXiv:2512.02927v2 (CC BY 4.0)}
\begin{document}
\maketitle

\noindent\emph{Source and licence.} Eisenstein cohomology and congruences for the ratios of Rankin--Selberg $L$-functions. Version arXiv:2512.02927v2 (CC BY 4.0), distributed by arXiv under the Creative Commons Attribution 4.0 International licence, http://creativecommons.org/licenses/by/4.0/, as recorded in the arXiv licence metadata for this exact version and shown on the abstract page https://arxiv.org/abs/2512.02927v2. Full licence notice: \texttt{licenses/arxiv-2512.02927-CC-BY-4.0.txt}.\par\smallskip

\noindent\textit{Editorial scope.} Counted unit: the article's lemma Lem:some-coset-reductions (source0.tex lines 1120--1130). The article's complete original proof of it is delivered with it (source0.tex lines 1132--1233, one \texttt{proof} environment of 102 lines). No mathematics is written, repaired or re-derived: the statement and the proof are the article's own lines, in the article's own notation, and no retained line is substituted or re-flowed.  No further source-local context is needed: every cross-reference of every retained line resolves inside the two delivered spans.  Nothing was omitted from any retained span, so the package carries no omission record.  No retained line contains a citation, so the wrapper carries no bibliography and no bibliography entry.  No retained line contains a figure environment, an image-inclusion command or drawing code, so no image is delivered and the package declares no asset dependency.  Editorial formatting only, in addition: an AMS \texttt{amsart} preamble that restates the article's own declarations and macros used by the retained lines, namely the environments \texttt{lemma} on separate counters, together with the standard packages those lines need.  The retained environments print in this document as Lemma 1 in order of appearance, whereas the article prints the same items with the numbering of its own full document.  The complete native source of the article as distributed in the arXiv e-print for this exact version is bundled unchanged as \texttt{sources/190227/source0.tex}.
\begin{lemma}
\label{Lem:some-coset-reductions}
\hfill%One has
\begin{enumerate}
\item[(i)] \begin{multline*}P(\mathbb{Q}_p)w_4\mirp{n_p + n'_p} = 	P(\mathbb{Q}_p)w_5\mirp{n_p + n'_p} \\ =
		P(\mathbb{Q}_p)w_6\mirp{n_p + n'_p} =	P(\mathbb{Q}_p)\xi_p^{(0)}\mirp{n_p + n'_p}.\end{multline*}

\smallskip
\item[(ii)] $K_p^M(\xi_p^{(0)}) = \kappa_P( P(\mathbb{Q}_p) \cap \xi_p^{(0)}  \mirp{n_p + n'_p} {\xi_p^{(0)}}^{-1} )  = K_p(n_p + n'_p) \times \GL_2(\mathbb{Z}_p).$
\end{enumerate}
\end{lemma}

\begin{proof}
For (i), observe that
$$%\begin{align*}
		w_4 \ = \ w_6\left(\begin{smallmatrix}
			0 & 1 & 0 & 0 \\
			0 & 0 & 1 & 0 \\
			1 & 0 & 0 & 0 \\
			0 & 0 & 0 & 1
		\end{smallmatrix}\right) 
		\quad \text{and} \quad 
		\left(\begin{smallmatrix}
			0 & 1 & 0 & 0 \\
			0 & 0 & 1 & 0 \\
			1 & 0 & 0 & 0 \\
			0 & 0 & 0 & 1
		\end{smallmatrix}\right) \in \mirp{n_p + n'_p}
$$
and, similarly, 
$$
w_5 \ = \ w_6 \left(\begin{smallmatrix}
			1 & 0 & 0 & 0 \\
			0 & 0 & 1 & 0 \\
			0 & 1 & 0 & 0 \\
			0 & 0 & 0 & 1
		\end{smallmatrix}\right)
		\quad \text{ and } \quad 
		\left(\begin{smallmatrix}
			1 & 0 & 0 & 0 \\
			0 & 0 & 1 & 0 \\
			0 & 1 & 0 & 0 \\
			0 & 0 & 0 & 1
		\end{smallmatrix}\right) \in \mirp{n_p + n'_p}.
$$
Hence $P(\mathbb{Q}_p)w_4 \mirp{n_p + n'_p} = P(\mathbb{Q}_p)w_5 \mirp{n_p + n'_p} =
	P(\mathbb{Q}_p)w_6 \mirp{n_p + n'_p}.$ To get the last equality of (i), further observe that
	\begin{equation*}
		w_6=
		\left(\begin{smallmatrix}
			1 & 0 & 0 & 0 \\
			0 & -1 & 0 & 1 \\
			0 & 0 & 1 & 0 \\
			0 & 0 & 0 & 1
		\end{smallmatrix}\right)
		\left(\begin{smallmatrix}
			1 & 0 & 0 & 0 \\
			0 & 1 & 0 & 0 \\
			0 & 0 & 1 & 0 \\
			0 & 1 & 0 & 1
		\end{smallmatrix}\right)
		\left(\begin{smallmatrix}
			1 & 0 & 0 & 0 \\
			0 & 1 & 0 & -1 \\
			0 & 0 & 1 & 0 \\
			0 & 0 & 0 & 1
		\end{smallmatrix}\right)
		\left(\begin{smallmatrix}
			0 & 0 & 1 & 0 \\
			0 & 1 & 0 & 0 \\
			1 & 0 & 0 & 0 \\
			0 & 0 & 0 & 1
		\end{smallmatrix}\right)
	\end{equation*}
	with
$%\begin{multline*}
		\left(\begin{smallmatrix}
			1 & 0 & 0 & 0 \\
			0 & -1 & 0 & 1 \\
			0 & 0 & 1 & 0 \\
			0 & 0 & 0 & 1
		\end{smallmatrix}\right) \in P(\mathbb{Q}_p)$ and 
$\left(\begin{smallmatrix}
			1 & 0 & 0 & 0 \\
			0 & 1 & 0 & -1 \\
			0 & 0 & 1 & 0 \\
			0 & 0 & 0 & 1
		\end{smallmatrix}\right), \left(\begin{smallmatrix}
			0 & 0 & 1 & 0 \\
			0 & 1 & 0 & 0 \\
			1 & 0 & 0 & 0 \\
			0 & 0 & 0 & 1
		\end{smallmatrix}\right) \in \mirp{n_p + n'_p}.$
This completes the proof of (i). 

For (ii), since 
$P(\mathbb{Q}_p) \cap \xi_p^{(0)}  \mirp{n_p + n'_p} {\xi_p^{(0)}}^{-1} = P(\mathbb{Q}_p) \cap w_6 \mirp{n_p + n'_p}w_6^{-1}$, 
and for an element 
$k =\left(\begin{smallmatrix}
		k_{11} & k_{12} & k_{13} & k_{14} \\
		k_{21} & k_{22} & k_{23} & k_{24} \\
		k_{31} & k_{32} & k_{33} & k_{34} \\
		k_{41} & k_{42} & k_{43} & k_{44}
	\end{smallmatrix}\right)$ in $\mirp{n_p + n'_p}$, 
one has
$w_6  k w_6^{-1} = 
	\left(\begin{smallmatrix}
		k_{33} & k_{34} & k_{31} & k_{32} \\
		k_{43} & k_{44} & k_{41} & k_{42} \\
		k_{13} & k_{14} & k_{11} & k_{12} \\
		k_{23} & k_{24} & k_{21} & k_{22}
	\end{smallmatrix}\right),$ 
(ii) follows. 
\end{proof}

\end{document}
